\chapter{The probability toolkit}\label{ch:toolkit}

\begin{chapterplan}
Which tools from probability and game theory does the paper use, and how does each tool work?
After this chapter you can:
\begin{itemize}
\item update a belief about the challenger's class with Bayes' rule, in its usual form and in odds form;
\item read a distribution function, its left limit and its atoms, and get expectations from it with the tail formula;
\item compare two distributions by first-order stochastic dominance, and say what a stronger incumbent is;
\item say what a best response, a mixed strategy, an equilibrium and a unilateral deviation are;
\item solve a linear threshold condition, and say when a strict inequality survives small changes of the parameters.
\end{itemize}
\end{chapterplan}

You can skim the tools that you already know. Later chapters cite each tool by its section, so you can come back
when a chapter needs one. Each section has the same order: an intuition in words, the statement, a full
derivation, a short illustration with the paper's objects, and one check-yourself question. No section uses
a numerical example.

The illustrations use four objects from the paper (Sections~2.2 and~4.2 of the paper):
\begin{itemize}
\item a challenger in one of two quality classes, \(H\) or \(L\), with the prior \(\pi=\Pr(j=H)\) and \(0<\pi<1\);
\item the distribution \(G_j\) of the challenger's acquisition value \(\theta\), given its class \(j\);
\item an incumbent whose value \(R\) is uniform on \([0,r]\), as in the paper's example;
\item a point mass at \(h\): a value that equals \(h\) with probability one, as for class \(H\) in the example.
\end{itemize}

The tools serve the later chapters in this way:
\begin{itemize}
\item Bayes' rule and its odds form (Sections~\ref{sec:toolkit-bayes} and~\ref{sec:toolkit-odds}) serve
Chapter~\ref{ch:beliefs}, where the price carries news about the class.
\item Distribution functions, expectations, the tail formula, maxima and dominance
(Sections~\ref{sec:toolkit-cdf} to~\ref{sec:toolkit-fosd}) serve the auction payoffs in Chapter~\ref{ch:auction}
and the example in Chapter~\ref{ch:example}.
\item The game-theory words (Section~\ref{sec:toolkit-games}) serve Chapters~\ref{ch:model}, \ref{ch:trading}
and~\ref{ch:coexistence}.
\item Monotone functions and thresholds (Section~\ref{sec:toolkit-monotone}) appear in every chapter from
Chapter~\ref{ch:beliefs} on.
\end{itemize}

% =====================================================================
\section{Events, conditional probability and Bayes' rule}\label{sec:toolkit-bayes}

\begin{intuition}[news moves a belief]
The challenger does not know its class. Suppose that it sees a signal, and that the signal is more likely in one
class than in the other. Bayes' rule turns the prior probability of class \(H\) into a posterior probability after
the signal. The signal counts in favour of the class that makes the signal more likely. The prior decides where
the belief starts.
\end{intuition}

\subsection{Events and conditional probability}

An event is a set of outcomes, such as ``the challenger is in class \(H\)'' or ``the order is a purchase''. For an
event \(A\), \(\Pr(A)\) is its probability, and ``not \(A\)'' has the probability \(1-\Pr(A)\). The event ``\(A\)
and \(B\)'' is written \(A\cap B\).

The conditional probability of \(A\) given \(B\) is
\begin{equation}\label{eq:toolkit-cond}
\Pr(A\mid B)=\frac{\Pr(A\cap B)}{\Pr(B)},\qquad \Pr(B)>0.
\end{equation}
In words: among the outcomes where \(B\) happens, it is the share where \(A\) also happens. Multiply both sides by
\(\Pr(B)\) to get the multiplication rule, \(\Pr(A\cap B)=\Pr(A\mid B)\Pr(B)\). The roles of \(A\) and \(B\) can
be swapped, so also \(\Pr(A\cap B)=\Pr(B\mid A)\Pr(A)\) when \(\Pr(A)>0\).

Two events are independent when
\[
\Pr(A\cap B)=\Pr(A)\Pr(B).
\]
When \(\Pr(B)>0\), divide by \(\Pr(B)\): independence says \(\Pr(A\mid B)=\Pr(A)\). News about \(B\) does not change
the probability of \(A\). Random variables are independent when every event about one is independent of every
event about the other. The paper uses this word often. For example, the preparation cost is independent of the
class and of the value, so the cost carries no information about either (Section~2.2 of the paper).

\subsection{The law of total probability}

Let a signal \(S\) take the value \(s\) with probability \(f_H(s)\) in class \(H\) and \(f_L(s)\) in class \(L\):
\[
f_H(s)=\Pr(S=s\mid H),\qquad f_L(s)=\Pr(S=s\mid L).
\]
These are probabilities of the signal given the class. They are not probabilities of the class. The two classes
split the outcomes: exactly one of them happens.

\begin{derivation}[the law of total probability]
\begin{align*}
\Pr(S=s)
&=\Pr(\{S=s\}\cap H)+\Pr(\{S=s\}\cap L) \why{\(H\) and \(L\) split the outcomes}\\
&=\Pr(S=s\mid H)\Pr(H)+\Pr(S=s\mid L)\Pr(L) \why{multiplication rule, twice}\\
&=\pi f_H(s)+(1-\pi)f_L(s) \why{notation, and \(\Pr(L)=1-\pi\)}
\end{align*}
\end{derivation}

The probability of the signal is an average of its two class probabilities. The weights are the class
probabilities.

\subsection{Bayes' rule}

Write \(\mu(s)=\Pr(H\mid S=s)\) for the posterior, the probability of class \(H\) after the signal \(s\).

\begin{derivation}[Bayes' rule]
Assume \(\Pr(S=s)>0\).
\begin{align}
\mu(s)=\Pr(H\mid S=s)
&=\frac{\Pr(H\cap\{S=s\})}{\Pr(S=s)} \why{definition \eqref{eq:toolkit-cond}}\notag\\
&=\frac{\Pr(S=s\mid H)\Pr(H)}{\Pr(S=s)} \why{multiplication rule, with the roles swapped}\notag\\
&=\frac{\pi f_H(s)}{\pi f_H(s)+(1-\pi)f_L(s)} \whyinline{law of total probability}\label{eq:toolkit-bayes}
\end{align}
\end{derivation}

The numerator is the probability of the path ``class \(H\), then signal \(s\)''. The denominator adds the
probabilities of the two paths that end in the signal \(s\). Figure~\ref{fig:toolkit-tree} shows the four paths.

\begin{figure}[tbp]
\centering
\begin{tikzpicture}[x=1cm,y=1cm,font=\small,
  lev/.style={inner sep=2pt}]
\node[lev] (root) at (0,0) {start};
\node[lev] (H) at (3.2,1.3) {class \(H\)};
\node[lev] (L) at (3.2,-1.3) {class \(L\)};
\node[lev,anchor=west] (Hs) at (6.4,2.0) {signal \(s\):\ \(\pi f_H(s)\)};
\node[lev,anchor=west] (Hn) at (6.4,0.6) {other signal:\ \(\pi[1-f_H(s)]\)};
\node[lev,anchor=west] (Ls) at (6.4,-0.6) {signal \(s\):\ \((1-\pi)f_L(s)\)};
\node[lev,anchor=west] (Ln) at (6.4,-2.0) {other signal:\ \((1-\pi)[1-f_L(s)]\)};
\draw[rulegrey] (root) -- node[above,sloped,text=ink] {\(\pi\)} (H);
\draw[rulegrey] (root) -- node[below,sloped,text=ink] {\(1-\pi\)} (L);
\draw[rulegrey] (H.east) -- node[above,sloped,text=ink] {\(f_H(s)\)} (Hs.west);
\draw[rulegrey] (H.east) -- (Hn.west);
\draw[rulegrey] (L.east) -- node[above,sloped,text=ink] {\(f_L(s)\)} (Ls.west);
\draw[rulegrey] (L.east) -- (Ln.west);
\end{tikzpicture}
\caption{A probability tree for the class and one signal. Each path probability is the product of the
probabilities along the path. The posterior after the signal \(s\) is the share of the top path among the two
paths that end in \(s\).}\label{fig:toolkit-tree}
\end{figure}

The posterior rises above the prior exactly when the signal is more likely in class
\(H\). To see this, compare \(\mu(s)\) with \(\pi\).
\begin{align*}
\mu(s)>\pi
&\iff\frac{f_H(s)}{\pi f_H(s)+(1-\pi)f_L(s)}>1 \why{divide \eqref{eq:toolkit-bayes} by \(\pi>0\)}\\
&\iff f_H(s)>\pi f_H(s)+(1-\pi)f_L(s) \why{multiply by the positive denominator}\\
&\iff(1-\pi)f_H(s)>(1-\pi)f_L(s) \why{subtract \(\pi f_H(s)\) from both sides}\\
&\iff f_H(s)>f_L(s) \why{divide by \(1-\pi>0\)}
\end{align*}
The same steps with ``\(<\)'' show that the posterior falls when \(f_H(s)<f_L(s)\). The same steps with
``\(=\)'' show that it equals the prior when \(f_H(s)=f_L(s)\). Look at two extreme cases. If
\(f_L(s)=0<f_H(s)\), the posterior is one: the signal proves the class. If \(\pi\) is near zero and both
\(f_H(s)\) and \(f_L(s)\) are positive, the posterior is also near zero. One signal of bounded strength cannot
overturn a strong prior.

In the paper, the signal is the order \(x\) that market makers see. The paper writes \(f_j(x)\) for the
probability of the order \(x\) in class \(j\), and \(\mu(x)\) for the posterior. Condition (E1) of
Definition~A.1 is formula~\eqref{eq:toolkit-bayes} with \(s=x\). Chapter~\ref{ch:beliefs} computes \(f_j(x)\) for
the paper's trading set-up.

\begin{watchout}
\begin{itemize}
\item Bayes' rule needs \(\Pr(S=s)>0\). After a signal of probability zero, the posterior is not defined by the
rule. In the paper every order has positive probability in both classes, so the rule applies after every order
(Lemma~A.1(a); Chapter~\ref{ch:beliefs}).
\item Keep \(\Pr(S=s\mid H)\) and \(\Pr(H\mid S=s)\) apart. The first is \(f_H(s)\), a property of the signal. The
second is the posterior \(\mu(s)\), which also depends on the prior.
\end{itemize}
\end{watchout}

\begin{checkyourself}\label{q:toolkit-bayes}
Use formula~\eqref{eq:toolkit-bayes}.
(a)~Suppose \(f_H(s)=f_L(s)>0\). Show that the posterior equals the prior.
(b)~Suppose \(f_L(s)=0<f_H(s)\). Find the posterior.
(c)~Fix \(f_H(s)>0\) and \(f_L(s)>0\). What happens to the posterior when the prior \(\pi\) moves toward one?
\end{checkyourself}

% =====================================================================
\section{Bayes' rule in odds form}\label{sec:toolkit-odds}

\begin{intuition}[one signal can move the odds only so far]
Odds compare two probabilities: the odds of class \(H\) are \(\Pr(H)/\Pr(L)\). In odds form, Bayes' rule is one
multiplication. The signal multiplies the prior odds by its likelihood ratio \(f_H(s)/f_L(s)\). Suppose that no
signal is more than \(1+u\) times as likely in one class as in the other. Then no signal can multiply the odds by
more than \(1+u\), or divide them by more than \(1+u\). So one signal can move the belief only inside a fixed
interval around the prior.
\end{intuition}

\subsection{The odds form}

The odds of a probability \(\mu\) with \(0<\mu<1\) are \(o=\mu/(1-\mu)\). The likelihood ratio of the signal \(s\)
is \(f_H(s)/f_L(s)\).

\begin{derivation}[posterior odds equal prior odds times the likelihood ratio]
Assume \(f_L(s)>0\). Start from Bayes' rule~\eqref{eq:toolkit-bayes} and write \(N=\pi f_H(s)+(1-\pi)f_L(s)\) for
its denominator.
\begin{align}
1-\mu(s)
&=\frac{N-\pi f_H(s)}{N} \why{write \(1\) as \(N/N\)}\notag\\
&=\frac{(1-\pi)f_L(s)}{N} \why{definition of \(N\)}\notag\\
\frac{\mu(s)}{1-\mu(s)}
&=\frac{\pi f_H(s)/N}{(1-\pi)f_L(s)/N} \why{divide the two formulas}\notag\\
&=\frac{\pi}{1-\pi}\cdot\frac{f_H(s)}{f_L(s)} \whyinline{\(N\) cancels; split the fraction}\label{eq:toolkit-odds}
\end{align}
\end{derivation}

This is the odds form of Bayes' rule in the paper's Appendix~A.2, with the order \(x\) in place of the signal \(s\).

To go back from odds to a probability, solve \(o=\mu/(1-\mu)\) for \(\mu\).
\begin{align*}
o(1-\mu)&=\mu \why{multiply by \(1-\mu>0\)}\\
o&=\mu+o\mu=\mu(1+o) \why{add \(o\mu\) to both sides}\\
\mu&=\frac{o}{1+o}=1-\frac1{1+o} \why{divide by \(1+o>0\)}
\end{align*}
As \(o\) rises, \(1/(1+o)\) falls, so \(\mu\) rises. The map from odds to probability is strictly increasing. The
rest of this section uses this fact: a bound on the odds gives a bound on the probability, in the same direction.

\subsection{The top and the bottom posterior}

\begin{keyidea}
Assume that every likelihood ratio lies between \(1/(1+u)\) and \(1+u\), for a number \(u>0\). Then every posterior
lies between the bottom posterior \(m(\pi)\) and the top posterior \(M(\pi)\):
\begin{equation}\label{eq:toolkit-bounds}
m(\pi)=\frac{\pi}{\pi+(1-\pi)(1+u)}\;\le\;\mu(s)\;\le\;M(\pi)=\frac{\pi(1+u)}{\pi(1+u)+1-\pi}.
\end{equation}
\end{keyidea}

In this chapter the bound on the likelihood ratio is an assumption. In the paper it is a result: the paper's
trading set-up gives the bound with \(u=\lambda/b\) (Section~2.3 of the paper). Chapter~\ref{ch:beliefs} derives
it, together with the condition on the trading set-up that it needs.

\begin{derivation}[the top posterior]
\begin{align*}
\frac{\mu(s)}{1-\mu(s)}
&=\frac{\pi}{1-\pi}\cdot\frac{f_H(s)}{f_L(s)} \why{odds form \eqref{eq:toolkit-odds}}\\
&\le\frac{\pi}{1-\pi}(1+u)=:o_{\max} \why{likelihood ratio at most \(1+u\); prior odds \(>0\)}
\end{align*}
The map \(o\mapsto o/(1+o)\) is increasing, so \(\mu(s)\le o_{\max}/(1+o_{\max})\). Compute this bound.
\begin{align*}
\frac{o_{\max}}{1+o_{\max}}
&=\frac{\pi(1+u)/(1-\pi)}{1+\pi(1+u)/(1-\pi)} \why{substitute \(o_{\max}\)}\\
&=\frac{\pi(1+u)/(1-\pi)}{[(1-\pi)+\pi(1+u)]/(1-\pi)} \why{common denominator below}\\
&=\frac{\pi(1+u)}{\pi(1+u)+1-\pi}=M(\pi) \why{cancel \(1-\pi\)}
\end{align*}
\end{derivation}

\begin{derivation}[the bottom posterior]
\begin{align*}
\frac{\mu(s)}{1-\mu(s)}
&=\frac{\pi}{1-\pi}\cdot\frac{f_H(s)}{f_L(s)} \why{odds form \eqref{eq:toolkit-odds}}\\
&\ge\frac{\pi}{(1-\pi)(1+u)}=:o_{\min} \why{likelihood ratio at least \(1/(1+u)\)}
\end{align*}
The map \(o\mapsto o/(1+o)\) is increasing, so \(\mu(s)\ge o_{\min}/(1+o_{\min})\).
\begin{align*}
\frac{o_{\min}}{1+o_{\min}}
&=\frac{\pi/[(1-\pi)(1+u)]}{1+\pi/[(1-\pi)(1+u)]} \why{substitute \(o_{\min}\)}\\
&=\frac{\pi/[(1-\pi)(1+u)]}{[(1-\pi)(1+u)+\pi]/[(1-\pi)(1+u)]} \why{common denominator below}\\
&=\frac{\pi}{\pi+(1-\pi)(1+u)}=m(\pi) \why{cancel \((1-\pi)(1+u)\)}
\end{align*}
\end{derivation}

A signal with the largest likelihood ratio \(1+u\) gives exactly \(M(\pi)\). A signal
with the smallest ratio \(1/(1+u)\) gives exactly \(m(\pi)\). The number \(u\) measures how strong one signal can
be. Test the two extremes.
\begin{itemize}
\item As \(u\) falls to zero, both \(1+u\) and \(1/(1+u)\) approach one. Then \(M(\pi)\) and \(m(\pi)\) both
approach \(\pi\): no signal can move the belief.
\item As \(u\) grows without limit, divide the numerator and the denominator of \(M(\pi)\) by \(1+u\):
\(M(\pi)=\pi/[\pi+(1-\pi)/(1+u)]\), which approaches \(\pi/\pi=1\). In \(m(\pi)\) the term \((1-\pi)(1+u)\) in the
denominator grows without limit, so \(m(\pi)\) approaches zero. A single signal can then come close to proof of the
class.
\end{itemize}

\subsection{Averages of posteriors}

Sometimes an observer does not see the signal itself. It sees only that the signal lies in a set \(A\) of values.
For example, the challenger sees a price, and two orders can have the same price. The observer's belief is
\(\Pr(H\mid S\in A)\). This belief is an average of the posteriors \(\mu(s)\) over the signals in \(A\).

\begin{derivation}[a coarse belief is an average of posteriors]
Assume \(\Pr(S\in A)>0\), and write \(w_s=\Pr(S=s)/\Pr(S\in A)\) for each \(s\in A\).
\begin{align*}
\Pr(H\mid S\in A)
&=\frac{\Pr(H\cap\{S\in A\})}{\Pr(S\in A)} \why{definition \eqref{eq:toolkit-cond}}\\
&=\frac{\sum_{s\in A}\Pr(H\cap\{S=s\})}{\Pr(S\in A)} \why{the events \(\{S=s\}\), \(s\in A\), do not overlap}\\
&=\frac{\sum_{s\in A}\Pr(H\mid S=s)\Pr(S=s)}{\Pr(S\in A)} \why{multiplication rule}\\
&=\sum_{s\in A}w_s\,\mu(s) \why{definitions of \(w_s\) and \(\mu(s)\)}
\end{align*}
Each weight \(w_s\) is nonnegative. The weights add up to one, because
\(\sum_{s\in A}\Pr(S=s)=\Pr(S\in A)\).
\end{derivation}

An average of numbers in an interval stays in the interval. Suppose \(m(\pi)\le\mu(s)\le M(\pi)\) for every \(s\in
A\).
\begin{align*}
\Pr(H\mid S\in A)=\sum_{s\in A}w_s\,\mu(s)
&\le\sum_{s\in A}w_s\,M(\pi) \why{\(\mu(s)\le M(\pi)\) and \(w_s\ge0\)}\\
&=M(\pi)\sum_{s\in A}w_s=M(\pi) \why{the weights add up to one}
\end{align*}
The same steps with \(\ge\) and \(m(\pi)\) give \(\Pr(H\mid S\in A)\ge m(\pi)\).

The paper's appendix uses a second route to the same end (Lemma~A.1, ``beliefs at a price''). In odds form, the
coarse belief has the likelihood ratio \(\sum_{s\in A}f_H(s)\big/\sum_{s\in A}f_L(s)\). A ratio of two sums of
positive terms lies between the smallest and the largest ratio of its terms. To see the upper half, suppose
\(f_H(s)\le K f_L(s)\) for each \(s\in A\). Add these inequalities over \(s\in A\) to get \(\sum f_H(s)\le K\sum
f_L(s)\), and divide by \(\sum f_L(s)>0\). The lower half is the same with \(\ge\).

\subsection{A symmetry of the two bounds}

The two bounds mirror each other:
\begin{equation}\label{eq:toolkit-mirror}
1-M(\pi)=m(1-\pi).
\end{equation}

\begin{derivation}[the mirror identity]
\begin{align*}
1-M(\pi)
&=\frac{\pi(1+u)+1-\pi-\pi(1+u)}{\pi(1+u)+1-\pi} \why{write \(1\) with the denominator of \(M(\pi)\)}\\
&=\frac{1-\pi}{\pi(1+u)+1-\pi} \why{cancel \(\pi(1+u)\) above}\\
m(1-\pi)
&=\frac{1-\pi}{(1-\pi)+[1-(1-\pi)](1+u)} \why{put \(1-\pi\) in place of \(\pi\) in \(m\)}\\
&=\frac{1-\pi}{(1-\pi)+\pi(1+u)} \why{\(1-(1-\pi)=\pi\)}
\end{align*}
The two right sides are equal.
\end{derivation}

In words: \(1-M(\pi)\) is the lowest probability of class \(L\) after one signal. Swap the names of the classes.
Class \(L\) then has the prior \(1-\pi\), and the bound on the likelihood ratio stays the same, because it is
symmetric. So the lowest posterior of class \(L\) is the bottom posterior at the prior \(1-\pi\).

The paper uses the mirror identity in its trading bound (Section~3.3 of the paper). There
\(\kappa(\pi)=\min\{m(\pi),1-M(\pi)\}\), and the identity turns it into \(\min\{m(\pi),m(1-\pi)\}\). The function
\(m\) is increasing: write \(m(\pi)=1/[1+(1+u)(1-\pi)/\pi]\), and note that \((1-\pi)/\pi\) falls as \(\pi\) rises.
So the smaller of \(m(\pi)\) and \(m(1-\pi)\) is \(m\) at the smaller of \(\pi\) and \(1-\pi\), and
\(\kappa(\pi)=m(\min\{\pi,1-\pi\})\), as in the paper.

\begin{watchout}
\begin{itemize}
\item The bounds are for one observation. The guide's own illustration, which the paper does not need: two
signals that are independent given the class multiply the odds twice. Their combined likelihood ratio can reach
\((1+u)^2\).
\item The average argument needs weights that are nonnegative and add up to one. The law of total probability
gives both.
\item The bound on the likelihood ratio must hold for every signal that the observer can see. Chapter~\ref{ch:beliefs}
shows which condition of the paper secures it for each order.
\end{itemize}
\end{watchout}

\begin{checkyourself}\label{q:toolkit-odds}
(a)~Show that \(m(\pi)<\pi<M(\pi)\) for every \(u>0\) and \(0<\pi<1\).
(b)~Show that \(M(1/2)=(1+u)/(2+u)\) and \(m(1/2)=1/(2+u)\), and that the two add up to one.
\end{checkyourself}

% =====================================================================
\section{Distribution functions, left limits and atoms}\label{sec:toolkit-cdf}

\begin{intuition}[the probability collected up to each level]
A distribution function answers one question at each level \(y\): how likely is the value to be at most \(y\)?
Read from left to right, the function collects probability. A jump shows a single value with positive
probability, an atom. A function without jumps belongs to a continuous distribution, where each single value has
probability zero.
\end{intuition}

\subsection{The distribution function}

The distribution function of a random value \(\theta\) is
\begin{equation}\label{eq:toolkit-cdf}
G(y)=\Pr(\theta\le y)\qquad\text{for every number }y.
\end{equation}
It has three properties.
\begin{itemize}
\item \(G\) is nondecreasing. If \(y<y'\), the event \(\{\theta\le y\}\) is part of the event \(\{\theta\le
y'\}\), so \(G(y)\le G(y')\).
\item \(G(y)\) is near zero for \(y\) far to the left, and near one for \(y\) far to the right.
\item \(G\) is continuous from the right. As \(y\) falls toward a level \(a\) from above, \(G(y)\) approaches
\(G(a)\).
\end{itemize}

The probability of an interval is a difference of two values of \(G\). For \(a<b\):
\begin{align*}
G(b)=\Pr(\theta\le b)
&=\Pr(\theta\le a)+\Pr(a<\theta\le b) \why{the two events do not overlap}\\
&=G(a)+\Pr(a<\theta\le b) \why{definition \eqref{eq:toolkit-cdf}}
\end{align*}
So \(\Pr(a<\theta\le b)=G(b)-G(a)\).

\subsection{The left limit and atoms}

The left limit \(G(p^-)\) is the number that \(G(y)\) approaches as \(y\) rises toward \(p\) from below. It is the
probability of a value strictly below \(p\):
\begin{equation}\label{eq:toolkit-left}
G(p^-)=\lim_{y\uparrow p}G(y)=\Pr(\theta<p).
\end{equation}
The reason in words: as \(y\) rises toward \(p\), the event \(\{\theta\le y\}\) grows and fills the event
\(\{\theta<p\}\). It never contains the value \(p\) itself.

An atom is a single value with positive probability. The jump of \(G\) at \(p\) is the probability of the value
\(p\).
\begin{align*}
\Pr(\theta=p)
&=\Pr(\theta\le p)-\Pr(\theta<p) \why{\(\{\theta\le p\}\) is \(\{\theta<p\}\) plus \(\{\theta=p\}\)}\\
&=G(p)-G(p^-) \why{definitions \eqref{eq:toolkit-cdf} and \eqref{eq:toolkit-left}}
\end{align*}
Two consequences follow from the complement rule:
\begin{equation}\label{eq:toolkit-geq}
\Pr(\theta\ge p)=1-G(p^-),\qquad \Pr(\theta>p)=1-G(p).
\end{equation}
A distribution is continuous when \(G\) has no jump, so that \(G(p^-)=G(p)\) at every \(p\). Then each single value
has probability zero, and \(\Pr(\theta\ge p)=\Pr(\theta>p)\).

\subsection{How the paper uses these objects}

\begin{itemize}
\item \emph{The reserve.} A bid equal to the reserve is admissible (Section~2.2 of the paper). A class-\(j\) value
fails the reserve when \(\theta<p\), and this has the probability \(G_j(p^-)\). The class ranking asks for
\(G_H(p^-)=G_L(p^-)=0\): every value meets the reserve.
\item \emph{Atoms in the classes.} The class distributions \(G_H\) and \(G_L\) may have atoms (Section~2.2 of the
paper). The example uses point masses, which are atoms of probability one.
\item \emph{The incumbent.} The incumbent's distribution \(F\) is continuous on \([0,\bar r]\), and it is
``extended by unity above \(\bar r\)'': \(F(y)=1\) for \(y>\bar r\). The incumbent's value never exceeds \(\bar r\),
so a value at most \(y\) is certain there. The extension lets a formula integrate \(F\) over \([p,\infty)\) with no
separate case above \(\bar r\). In the same way, \(F(y)=0\) for \(y<0\), because \(R\ge0\).
\end{itemize}

\subsection{Two illustrations}

A point mass at \(h\) has the distribution function \(G(y)=\1\{y\ge h\}\). Here \(\1\{\cdot\}\) is the
indicator: it equals one when the condition in the braces holds and zero otherwise. Check the two cases.
\begin{itemize}
\item If \(y<h\), the value \(h\) lies above \(y\), so the event \(\{\theta\le y\}\) never happens: \(G(y)=0\).
\item If \(y\ge h\), the value \(h\) is at most \(y\), so the event always happens: \(G(y)=1\).
\end{itemize}
So \(G(h^-)=0\) and \(G(h)=1\). The jump at \(h\) has height one: an atom of probability one. By
\eqref{eq:toolkit-geq}, \(\Pr(\theta\ge h)=1-0=1\) and \(\Pr(\theta>h)=1-1=0\).

The uniform distribution on \([0,r]\) spreads its probability evenly over \([0,r]\), with the density \(1/r\). Its
distribution function is
\begin{equation}\label{eq:toolkit-uniform}
F(y)=\begin{cases}0,& y<0,\\ y/r,& 0\le y\le r,\\ 1,& y>r.\end{cases}
\end{equation}
For the middle line, add up the density from \(0\) to \(y\):
\begin{align*}
F(y)=\Pr(R\le y)&=\int_0^y\frac1r\,dz \why{the density is \(1/r\) on \([0,r]\)}\\
&=\frac yr \why{a constant integrand times the length \(y\)}
\end{align*}
At \(y=r\) the middle line gives \(r/r=1\), the value above \(r\). At \(y=0\) it gives \(0\), the value below
zero. So \(F\) has no jump: the uniform law is continuous, and \(\Pr(R=y)=0\) for every \(y\). Here
\(\bar r=r\), and for \(y\ge0\) the three lines combine into \(F(y)=\min\{y/r,1\}\). Figure~\ref{fig:toolkit-cdfs}
draws both functions.

\begin{figure}[tbp]
\centering
\begin{tikzpicture}[x=1cm,y=2.2cm,font=\small]
\draw[-{Latex[length=2mm]}] (0,0) -- (6.2,0) node[below] {\(y\)};
\draw[-{Latex[length=2mm]}] (0,0) -- (0,1.3);
\draw (0.08,1) -- (-0.08,1) node[left] {\(1\)};
\node[left] at (0,0) {\(0\)};
\draw[dashed,rulegrey] (3.5,0) -- (3.5,1);
\draw[accent,very thick] (0,0) -- (3.5,0);
\draw[accent,very thick] (3.5,1) -- (5.9,1);
\filldraw[fill=white,draw=accent,thick] (3.5,0) circle (2pt);
\fill[accent] (3.5,1) circle (2pt);
\node[below] at (3.5,0) {\(h\)};
\node[above] at (4.7,1) {\(G(y)\)};
\node at (3.1,-0.42) {(a) a point mass at \(h\)};
\begin{scope}[xshift=8.2cm]
\draw[-{Latex[length=2mm]}] (0,0) -- (6.2,0) node[below] {\(y\)};
\draw[-{Latex[length=2mm]}] (0,0) -- (0,1.3);
\draw (0.08,1) -- (-0.08,1) node[left] {\(1\)};
\node[left] at (0,0) {\(0\)};
\draw[dashed,rulegrey] (4,0) -- (4,1);
\draw[accent,very thick] (0,0) -- (4,1) -- (5.9,1);
\node[below] at (4,0) {\(r\)};
\node[above left] at (2.3,0.55) {\(F(y)=y/r\)};
\node at (3.1,-0.42) {(b) uniform on \([0,r]\)};
\end{scope}
\end{tikzpicture}
\caption{Two distribution functions. In (a) the function jumps from zero to one at \(h\): the filled dot shows
that \(G(h)=1\), and the open dot shows that \(G(h^-)=0\). In (b) the function rises in a straight line and has
no jump, so the distribution is continuous.}\label{fig:toolkit-cdfs}
\end{figure}

\begin{watchout}
\begin{itemize}
\item \(G(p^-)\) and \(G(p)\) differ by the atom at \(p\). For a point mass at the reserve, \(G(p^-)=0\) but
\(G(p)=1\). A mix-up turns ``every value meets the reserve'' into ``no value passes the reserve''.
\item The paper treats the two kinds of distribution differently. The class distributions \(G_j\) may have atoms.
The incumbent's distribution \(F\) is continuous, so that a tie \(R=\theta\) has probability zero
(Section~2.2 of the paper; Section~\ref{sec:toolkit-maxmin} below).
\end{itemize}
\end{watchout}

\begin{checkyourself}\label{q:toolkit-cdf}
Let class \(L\) be a point mass at \(\ell\), and suppose that the reserve equals \(\ell\), so \(p=\ell\).
(a)~Find \(G_L(p^-)\) and \(G_L(p)\).
(b)~Does every class-\(L\) value meet the reserve?
(c)~Which of the two numbers does the class ranking restrict, and does the restriction hold here?
\end{checkyourself}

% =====================================================================
\section{Expectation, conditional expectation and averages over classes}\label{sec:toolkit-expectation}

\begin{intuition}[average in two stages]
An expectation is an average with probabilities as weights. Two facts do most of the work in the paper. First,
the average of a sum is the sum of the averages. Second, you can average in two stages: first inside each class,
then across the classes, with the class probabilities as weights. The second fact makes every expectation at a
belief \(\mu\) a straight line in \(\mu\).
\end{intuition}

\subsection{Expectation and linearity}

If \(X\) takes finitely many values, its expectation is \(\E[X]=\sum_x x\Pr(X=x)\). For a function \(\varphi\),
\(\E[\varphi(X)]=\sum_x\varphi(x)\Pr(X=x)\). The illustrations of this chapter need two special cases.
\begin{itemize}
\item For a point mass at \(h\), \(\E[\varphi(\theta)]=\varphi(h)\). There is one value, with probability one.
\item For the uniform law on \([0,r]\), \(\E[\varphi(R)]=\frac1r\int_0^r\varphi(y)\,dy\). The integral replaces
the sum, and the density \(1/r\) replaces the probability of each value. With \(\varphi(y)=y\), this gives
\(\E[R]=\frac1r\cdot\frac{r^2}2=\frac r2\).
\end{itemize}

Expectation is linear: for numbers \(a\), \(b\), \(c\),
\begin{equation}\label{eq:toolkit-linear}
\E[aX+bY+c]=a\E[X]+b\E[Y]+c.
\end{equation}

\begin{derivation}[linearity]
Take a finite set of outcomes \(\omega\), each with the probability \(\Pr(\omega)\). The variables \(X\) and \(Y\)
take the values \(X(\omega)\) and \(Y(\omega)\).
\begin{align*}
\E[aX+bY+c]
&=\sum_\omega[aX(\omega)+bY(\omega)+c]\Pr(\omega) \why{definition, outcome by outcome}\\
&=a\sum_\omega X(\omega)\Pr(\omega)+b\sum_\omega Y(\omega)\Pr(\omega) \why{split the sum}\\
&\qquad+c\sum_\omega\Pr(\omega)\\
&=a\E[X]+b\E[Y]+c \why{definition; probabilities add up to one}
\end{align*}
\end{derivation}

Linearity needs no independence between \(X\) and \(Y\).

The expectation of an indicator is a probability:
\begin{equation}\label{eq:toolkit-indicator}
\E\bigl[\1\{A\}\bigr]=1\cdot\Pr(A)+0\cdot\bigl(1-\Pr(A)\bigr)=\Pr(A).
\end{equation}
As an illustration, take the target's expected proceeds without entry, \(t_0=p[1-F(p)]\) (Section~2.2 of the
paper). Without a challenger, the incumbent buys at the reserve when its value meets the reserve. So the
proceeds are \(p\,\1\{R\ge p\}\).
\begin{align*}
t_0=\E\bigl[p\,\1\{R\ge p\}\bigr]
&=p\,\E\bigl[\1\{R\ge p\}\bigr] \why{linearity \eqref{eq:toolkit-linear}}\\
&=p\Pr(R\ge p) \why{expectation of an indicator \eqref{eq:toolkit-indicator}}\\
&=p[1-F(p^-)] \why{complement rule \eqref{eq:toolkit-geq}}\\
&=p[1-F(p)] \why{\(F\) is continuous}\\
&=p\Bigl(1-\frac pr\Bigr) \why{uniform \eqref{eq:toolkit-uniform}, with \(p\le r\)}
\end{align*}
The last line is the paper's \(t_0\) in the example, where \(p<r\) (Appendix~A.1).

\subsection{Conditional expectation given the class}

The class expectation \(\E_j[\varphi(\theta)]\) is the expectation of \(\varphi(\theta)\) given the class \(j\). It
uses the class distribution \(G_j\). The paper defines its class objects in this way (Section~2.4 of the paper):
\[
t_j=\E_j[t(\theta)],\qquad g_j=\E_j[g(\theta)],\qquad j\in\{H,L\}.
\]
Here \(t(\theta)\) is the target's expected proceeds after the challenger prepares, at a known challenger value
\(\theta\). In the same way, \(g(\theta)\) is the challenger's expected gross acquisition profit at the value
\(\theta\). In the example the
classes are point masses at \(h\) and \(\ell\). So the class expectation is the function at one point:
\(t_H=t(h)\), \(t_L=t(\ell)\), \(g_H=g(h)\) and \(g_L=g(\ell)\).

\subsection{The law of iterated expectations}

The overall expectation is the average of the class expectations, with the class probabilities as weights:
\begin{equation}\label{eq:toolkit-lie}
\E[\varphi(\theta)]=\pi\,\E_H[\varphi(\theta)]+(1-\pi)\,\E_L[\varphi(\theta)].
\end{equation}

\begin{derivation}[an average over the classes]
Take a value \(\theta\) with finitely many values \(y\).
\begin{align*}
\E[\varphi(\theta)]
&=\sum_y\varphi(y)\Pr(\theta=y) \why{definition}\\
&=\sum_y\varphi(y)\bigl[\pi\Pr(\theta=y\mid H)+(1-\pi)\Pr(\theta=y\mid L)\bigr]
\why{total probability (Section~\ref{sec:toolkit-bayes})}\\
&=\pi\sum_y\varphi(y)\Pr(\theta=y\mid H) \why{split the sum}\\
&\qquad+(1-\pi)\sum_y\varphi(y)\Pr(\theta=y\mid L)\\
&=\pi\,\E_H[\varphi(\theta)]+(1-\pi)\,\E_L[\varphi(\theta)] \why{definition of \(\E_j\)}
\end{align*}
In the third line, \(\pi\) and \(1-\pi\) move out of the sums because they do not depend on \(y\). With a density,
the sums become integrals and the steps are the same.
\end{derivation}

\subsection{Expectations at a belief}

Suppose that an observer sees a signal \(S=s\) and holds the belief \(\mu=\Pr(H\mid S=s)\). Assume that the signal
is independent of the value given the class:
\[
\Pr(\theta=y\mid j,\ S=s)=\Pr(\theta=y\mid j)\qquad\text{for each class }j\text{ and each value }y.
\]

\begin{derivation}[an expectation at a belief]
\begin{align*}
\E[\varphi(\theta)\mid S=s]
&=\sum_y\varphi(y)\Pr(\theta=y\mid S=s) \why{definition}\\
&=\sum_y\varphi(y)\bigl[\mu\Pr(\theta=y\mid H,S=s) \why{total probability, given \(S=s\)}\\
&\qquad+(1-\mu)\Pr(\theta=y\mid L,S=s)\bigr]\\
&=\sum_y\varphi(y)\bigl[\mu\Pr(\theta=y\mid H)+(1-\mu)\Pr(\theta=y\mid L)\bigr]
\why{independent given the class}\\
&=\mu\,\E_H[\varphi(\theta)]+(1-\mu)\,\E_L[\varphi(\theta)] \why{split the sum}
\end{align*}
\end{derivation}

The signal changes only the weights on the two classes. It does not change the class distributions \(G_H\) and
\(G_L\). Apply this to the gross profit, \(\varphi=g\):
\begin{align}
B(\mu)&=\mu\,g_H+(1-\mu)\,g_L \why{the formula above, with \(\varphi=g\)}\notag\\
&=g_L+\mu\,(g_H-g_L) \whyinline{collect the terms in \(\mu\)}\label{eq:toolkit-B}
\end{align}
This is the paper's formula for gross profit at a belief (Section~2.4 of the paper). The paper gives the reason in
one phrase: the challenger's value depends on the price only through its class.

In words, \(B\) is a straight line in \(\mu\). It equals \(g_L\) at \(\mu=0\) and \(g_H\) at \(\mu=1\), and its
slope \(g_H-g_L\) is the profit gap between the classes. A belief halfway between zero and one gives a profit
halfway between \(g_L\) and \(g_H\). At the extreme \(g_H=g_L\), the line is flat, and news about the class does not
change gross profit.

The same argument makes the expected proceeds at a belief a straight line, at a fixed probability of entry. The
paper's price after an order \(x\) has this form (Section~2.4 of the paper):
\[
P(x)=t_0+e(x)\,\bigl[t_L-t_0+\DT\,\mu(x)\bigr].
\]
Here \(e(x)\) is entry after the order, and \(\DT=t_H-t_L\) is the target-payoff spread. Later chapters derive
this formula.

\begin{watchout}
\begin{itemize}
\item \(B(\mu)\) is a straight line because the signal is independent of the value given the class. Without this
independence, the signal would also change the class distributions, and \(B\) would not be linear in \(\mu\). The
paper's extension with private signals of the investor and the challenger keeps this hypothesis (Chapter~\ref{ch:robustness}).
\item The price is a straight line in the posterior only at a fixed level of entry. In equilibrium, entry also
depends on the posterior, so equilibrium prices also reflect entry.
\end{itemize}
\end{watchout}

\begin{checkyourself}\label{q:toolkit-expectation}
The paper writes \(w_j=t_j-t_0\) for the gain in expected proceeds that entry brings in class \(j\), and
\(w(\mu)=(1-\mu)w_L+\mu w_H\) for its average at the belief \(\mu\) (Appendix~A.2).
(a)~Show that \(w(\mu)=t_L-t_0+\DT\,\mu\).
(b)~What are \(w(0)\), \(w(1)\) and the slope of \(w\)?
\end{checkyourself}

% =====================================================================
\section{The tail formula}\label{sec:toolkit-tail}

\begin{intuition}[an expectation adds up thin slices]
Think of a nonnegative value \(X\) as a stack of thin horizontal slices, from the floor at zero up to the height
\(X\). The slice at the height \(y\) is in the stack exactly when \(X\) passes \(y\). To find the average height of
the stack, add up all the slices, and weight each slice by the probability that it is present. That probability
is \(\Pr(X>y)=1-G(y)\). So the expectation is the area under the curve \(1-G(y)\).
\end{intuition}

\subsection{The formula}

\begin{keyidea}
For a random value \(X\ge0\) with the distribution function \(G\),
\begin{equation}\label{eq:toolkit-tail}
\E[X]=\int_0^\infty\Pr(X>y)\,dy=\int_0^\infty[1-G(y)]\,dy.
\end{equation}
\end{keyidea}

The integral adds up, over every level \(y\) from zero upward, the probability that \(X\) passes the level.

\begin{derivation}[the tail formula]
First write \(X\) itself as an integral. The length of the interval \([0,X]\) is \(X\):
\begin{align*}
X&=\int_0^X 1\,dy \why{the length of \([0,X]\)}\\
&=\int_0^\infty\1\{y<X\}\,dy \why{the indicator is one below \(X\) and zero from \(X\) on}
\end{align*}
Then take the expectation of both sides.
\begin{align*}
\E[X]&=\E\Bigl[\int_0^\infty\1\{y<X\}\,dy\Bigr] \why{take expectations}\\
&=\int_0^\infty\E\bigl[\1\{y<X\}\bigr]\,dy \why{swap the expectation and the integral (see below)}\\
&=\int_0^\infty\Pr(X>y)\,dy \why{expectation of an indicator \eqref{eq:toolkit-indicator}}\\
&=\int_0^\infty[1-G(y)]\,dy \why{complement rule \eqref{eq:toolkit-geq}}
\end{align*}
\end{derivation}

The second step swaps two operations: an average over outcomes and a sum over levels. The analogy below explains
why the swap is allowed.

\begin{analogy}[a table added up in two orders]
Picture a table. Each row is an outcome \(\omega\), and each column is a thin level from \(y\) to \(y+dy\). The
entry in row \(\omega\) and column \(y\) is \(\Pr(\omega)\cdot\1\{y<X(\omega)\}\,dy\). Add each row along the
columns first: row \(\omega\) gives \(\Pr(\omega)X(\omega)\), and the rows together give \(\E[X]\). Add each column
over the rows first: column \(y\) gives \(\Pr(X>y)\,dy\), and the columns together give the integral. Both orders
add up the same entries, so they give the same total.

Where the analogy breaks: with infinitely many rows or
columns, entries of both signs can give two different totals in the two orders. Nonnegative entries always give
one total. The paper cites Tonelli's theorem for this step (Appendix~A.1), and it is the reason for \(X\ge0\).
\end{analogy}

A remark on ``\(>\)'' and ``\(\ge\)''. The two probabilities \(\Pr(X>y)\) and \(\Pr(X\ge y)\) differ by
\(\Pr(X=y)\), which is positive only at atoms. A distribution has at most one atom of probability above
\(\tfrac12\), at most two above \(\tfrac13\), and so on, so its atoms form a list. An integral does not change when
the integrand changes at a list of single points. So \(\E[X]=\int_0^\infty\Pr(X\ge y)\,dy\) as well. The paper uses
this remark for the class distributions (Appendix~A.1).

Figure~\ref{fig:toolkit-tail}(a) shows the formula as an area. The thin strip at the level \(y\) has the width
\(dy\) and the height \(1-G(y)\).

\begin{figure}[tbp]
\centering
\begin{tikzpicture}[x=1cm,y=2.4cm,font=\small]
\fill[accent!15] (0,0) -- plot[domain=0:5.6,samples=60] (\x,{exp(-0.55*\x)}) -- (5.6,0) -- cycle;
\fill[accent!55] (1.6,0) rectangle (1.8,{exp(-0.55*1.7)});
\draw[accent,very thick] plot[domain=0:5.6,samples=60] (\x,{exp(-0.55*\x)});
\draw[-{Latex[length=2mm]}] (0,0) -- (6.2,0) node[below] {\(y\)};
\draw[-{Latex[length=2mm]}] (0,0) -- (0,1.25);
\draw (0.08,1) -- (-0.08,1) node[left] {\(1\)};
\node[left] at (0,0) {\(0\)};
\node[below] at (1.7,0) {\(y\)};
\draw[->,notegrey] (2.6,0.75) node[right,text=ink] {\(1-G(y)\)} -- (1.85,0.45);
\node at (3.1,-0.4) {(a) \(\E[X]\) is the whole area};
\begin{scope}[xshift=8.2cm]
\fill[accent!15] (0,0) -- plot[domain=0:2.2,samples=30] (\x,{exp(-0.55*\x)}) -- (2.2,0) -- cycle;
\fill[accent!45] (2.2,0) -- plot[domain=2.2:5.6,samples=40] (\x,{exp(-0.55*\x)}) -- (5.6,0) -- cycle;
\draw[accent,very thick] plot[domain=0:5.6,samples=60] (\x,{exp(-0.55*\x)});
\draw[-{Latex[length=2mm]}] (0,0) -- (6.2,0) node[below] {\(y\)};
\draw[-{Latex[length=2mm]}] (0,0) -- (0,1.25);
\draw (0.08,1) -- (-0.08,1) node[left] {\(1\)};
\node[left] at (0,0) {\(0\)};
\draw[dashed,rulegrey] (2.2,0) -- (2.2,{exp(-0.55*2.2)});
\node[below] at (2.2,0) {\(c\)};
\node[font=\footnotesize] at (0.95,0.2) {\(\E[\min\{X,c\}]\)};
\draw[->,notegrey] (4.1,0.62) node[above,text=ink] {\(\E[(X-c)_+]\)} -- (3.1,0.1);
\node at (3.1,-0.4) {(b) the area split at \(c\)};
\end{scope}
\end{tikzpicture}
\caption{The tail formula as an area. The curve is \(1-G(y)\), the probability that \(X\) passes the level \(y\).
In (a), the whole shaded area is \(\E[X]\), and the dark strip is one slice. In (b), the area to the left of \(c\)
is \(\E[\min\{X,c\}]\), and the area to the right of \(c\) is \(\E[(X-c)_+]\), where \((x)_+=\max\{x,0\}\).}\label{fig:toolkit-tail}
\end{figure}

Two illustrations confirm the formula against the direct computation.
\begin{itemize}
\item A point mass at \(h\) has \(1-G(y)=\1\{y<h\}\). The formula gives \(\int_0^\infty\1\{y<h\}\,dy=h\), the
value itself.
\item The uniform law on \([0,r]\) has \(1-F(y)=1-y/r\) on \([0,r]\) and zero above \(r\). The formula gives
\[
\int_0^r\Bigl(1-\frac yr\Bigr)dy=\Bigl[y-\frac{y^2}{2r}\Bigr]_0^r=r-\frac r2=\frac r2,
\]
the mean \(\E[R]=r/2\) of Section~\ref{sec:toolkit-expectation}.
\end{itemize}

In words: a distribution that passes every level with a larger probability has a larger mean. Section~\ref{sec:toolkit-fosd} builds on this point.

\subsection{Four formulas that follow}

Four versions of the tail formula serve the later chapters. Each one cuts the stack of slices at a level. The
paper writes \((x)_+=\max\{x,0\}\) for the positive part of a number \(x\), and this guide keeps that notation.
Section~\ref{sec:toolkit-monotone} lists its rules.

\begin{keyidea}
Let \(X\ge0\) have the distribution function \(G\), and let \(c\ge0\) and \(p\ge0\). Then
\begin{align}
\E[(X-c)_+]&=\int_c^\infty[1-G(y)]\,dy,\label{eq:toolkit-pospart}\\
\E[\max\{X,p\}]&=p+\int_p^\infty[1-G(y)]\,dy,\label{eq:toolkit-max}\\
\E[\min\{X,c\}]&=\int_0^c[1-G(y)]\,dy.\label{eq:toolkit-min}
\end{align}
Let \(R\) have a continuous distribution function \(F\), and fix a number \(\theta\ge p\). Then
\begin{equation}\label{eq:toolkit-fixed}
\E_R\bigl[(\theta-\max\{p,R\})_+\bigr]=\int_p^\theta F(y)\,dy.
\end{equation}
\end{keyidea}

Each proof has the same two steps. First, write the variable as an integral of an indicator, and check the
cases. Second, take expectations and swap, as in the tail formula.

\begin{derivation}[the slices above \(c\): formula \eqref{eq:toolkit-pospart}]
The claim is \((X-c)_+=\int_c^\infty\1\{y<X\}\,dy\). Check the two cases.
\begin{itemize}
\item If \(X\le c\), the left side is zero. On the right, every \(y\ge c\) has \(y\ge X\), so the indicator is
zero and the integral is zero.
\item If \(X>c\), the left side is \(X-c\). On the right, the indicator is one for \(c\le y<X\) and zero for
\(y\ge X\), so the integral is the length \(X-c\).
\end{itemize}
Take expectations:
\begin{align*}
\E[(X-c)_+]&=\int_c^\infty\E\bigl[\1\{y<X\}\bigr]\,dy \why{swap; the integrand is nonnegative}\\
&=\int_c^\infty[1-G(y)]\,dy \why{\(\E[\1\{y<X\}]=\Pr(X>y)=1-G(y)\)}
\end{align*}
\end{derivation}

\begin{derivation}[a floor at \(p\): formula \eqref{eq:toolkit-max}]
First, \(\max\{X,p\}=p+(X-p)_+\). If \(X\ge p\), both sides equal \(X\). If \(X<p\), both sides equal \(p\). Then
\begin{align*}
\E[\max\{X,p\}]&=p+\E[(X-p)_+] \why{linearity \eqref{eq:toolkit-linear}}\\
&=p+\int_p^\infty[1-G(y)]\,dy \why{formula \eqref{eq:toolkit-pospart} with \(c=p\)}
\end{align*}
\end{derivation}

In words: a floor at \(p\) gives \(p\) for sure, plus the slices above \(p\). The paper uses this formula for
\(\E_j[\max\{\theta,p\}]\) (Appendix~A.1).

\begin{derivation}[the slices below \(c\): formula \eqref{eq:toolkit-min}]
The claim is \(\min\{X,c\}=\int_0^c\1\{y<X\}\,dy\). Check the two cases.
\begin{itemize}
\item If \(X\ge c\), the indicator is one for every \(y\) in \([0,c)\), so the integral is \(c=\min\{X,c\}\).
\item If \(X<c\), the indicator is one for \(0\le y<X\) and zero for \(X\le y<c\), so the integral is
\(X=\min\{X,c\}\).
\end{itemize}
Take expectations:
\begin{align*}
\E[\min\{X,c\}]&=\int_0^c\E\bigl[\1\{y<X\}\bigr]\,dy \why{swap; the integrand is nonnegative}\\
&=\int_0^c[1-G(y)]\,dy \why{expectation of an indicator}
\end{align*}
\end{derivation}

The last two cuts fit together. In every case \(\min\{X,c\}+(X-c)_+=X\): the slices below \(c\) plus the slices
above \(c\) make the whole stack. In the integrals, \(\int_0^c+\int_c^\infty=\int_0^\infty\).
Figure~\ref{fig:toolkit-tail}(b) shows the split.

\begin{derivation}[a fixed value against a random incumbent: formula \eqref{eq:toolkit-fixed}]
Fix \(\theta\ge p\). The claim is \((\theta-\max\{p,R\})_+=\int_p^\theta\1\{R<y\}\,dy\). Check the three cases.
\begin{itemize}
\item If \(R<p\), then \(\max\{p,R\}=p\), and the left side is \(\theta-p\). On the right, every \(y\ge p\) has
\(R<y\), so the integral is the length \(\theta-p\).
\item If \(p\le R<\theta\), then \(\max\{p,R\}=R\), and the left side is \(\theta-R\). On the right, the indicator
is zero for \(p\le y\le R\) and one for \(R<y\le\theta\), so the integral is \(\theta-R\).
\item If \(R\ge\theta\), the left side is zero. On the right, every \(y\le\theta\) has \(y\le R\), so the
indicator is zero and the integral is zero.
\end{itemize}
Take the expectation over \(R\):
\begin{align*}
\E_R\bigl[(\theta-\max\{p,R\})_+\bigr]
&=\int_p^\theta\Pr(R<y)\,dy \why{swap; expectation of an indicator}\\
&=\int_p^\theta F(y)\,dy \why{\(F\) is continuous, so \(\Pr(R<y)=F(y)\)}
\end{align*}
\end{derivation}

The paper reads this formula as an acquisition profit. A prepared challenger with the value \(\theta\ge p\)
keeps \((\theta-\max\{p,R\})_+\) (Section~3.1 of the paper). So its expected gross profit is
\(g(\theta)=\int_p^\theta F(y)\,dy\) (Appendix~A.1). The slices picture explains it. Cut the challenger's
value between \(p\) and \(\theta\) into thin slices. The challenger keeps the slice at the level \(y\) exactly when
the incumbent's value lies below \(y\), and that has the probability \(F(y)\).

As an illustration, take the uniform incumbent on \([0,r]\) and a point mass at \(h\), with \(p<r<h\). The integral
adds up \(F(y)\) from \(p\) to \(h\), and \(F\) has a different formula below and above \(r\).
\begin{align*}
\int_p^hF(y)\,dy
&=\int_p^r\frac yr\,dy+\int_r^h1\,dy \why{\(F(y)=y/r\) below \(r\), and \(F=1\) above \(r\)}\\
&=\frac{r^2-p^2}{2r}+(h-r) \why{antiderivative \(y^2/(2r)\); length of \([r,h]\)}\\
&=h-\frac r2-\frac{p^2}{2r} \why{\(r^2/(2r)=r/2\); collect the terms}
\end{align*}
This is the paper's \(g_H\) in the example (Section~4.2 of the paper). The second integral uses the extension of
\(F\) by unity above \(r\). Between \(r\) and \(h\), the challenger keeps every slice, because the incumbent never
reaches those levels.

\begin{watchout}
\begin{itemize}
\item The tail formula in this form needs \(X\ge0\). The nonnegative integrand is what permits the swap.
\item The formulas for \(\max\{X,p\}\) and \(\min\{X,c\}\) need \(p\ge0\) and \(c\ge0\), because the slices start at
zero. In the paper, \(p>0\).
\item In formulas \eqref{eq:toolkit-pospart} to \eqref{eq:toolkit-min} the integrand is the tail \(1-G\). In
formula \eqref{eq:toolkit-fixed} it is \(F\) itself, because the challenger keeps a slice when the incumbent lies
below it.
\end{itemize}
\end{watchout}

\begin{checkyourself}\label{q:toolkit-tail}
Take the uniform incumbent on \([0,r]\) and a point mass at \(\ell\), with \(p<\ell<r\).
(a)~Use formula~\eqref{eq:toolkit-fixed} to find \(g(\ell)=\int_p^\ell F(y)\,dy\).
(b)~Say how \(g(\ell)\) changes when \(r\) rises, and what happens as \(\ell\) approaches \(p\).
\end{checkyourself}

% =====================================================================
\section{Maxima, minima and the probability to win}\label{sec:toolkit-maxmin}

\begin{intuition}[``both'' becomes a product]
The larger of two values is at most \(y\) exactly when both values are at most \(y\). The smaller of two values is
above \(y\) exactly when both values are above \(y\). When the two values are independent, the probability of
``both'' is the product of the two probabilities. A fixed value beats a random value when the random value falls
below it, and with a continuous distribution a tie has probability zero.
\end{intuition}

\subsection{The maximum and the minimum of two independent values}

Let \(X\) and \(Y\) be independent, with the distribution functions \(F_X\) and \(F_Y\).

\begin{derivation}[the maximum and the minimum]
\begin{align}
\Pr(\max\{X,Y\}\le y)
&=\Pr(X\le y\text{ and }Y\le y) \why{the larger is at most \(y\) when both are}\notag\\
&=\Pr(X\le y)\Pr(Y\le y) \why{independence}\notag\\
&=F_X(y)\,F_Y(y) \whyinline{definition \eqref{eq:toolkit-cdf}}\label{eq:toolkit-maxcdf}\\[4pt]
\Pr(\min\{X,Y\}>y)
&=\Pr(X>y\text{ and }Y>y) \why{the smaller is above \(y\) when both are}\notag\\
&=\Pr(X>y)\Pr(Y>y) \why{independence}\notag\\
&=[1-F_X(y)]\,[1-F_Y(y)] \whyinline{complement rule \eqref{eq:toolkit-geq}}\label{eq:toolkit-mintail}
\end{align}
\end{derivation}

In words: the maximum stays at or below a level only if both values do, and the minimum passes a level only if
both values pass it. Each formula multiplies two numbers between zero and one, so the product is at most
each factor.

\subsection{A fixed value against a random value}

Fix a number \(\theta\), and let \(R\) have a continuous distribution function \(F\).
\begin{align*}
\Pr(R<\theta)&=F(\theta^-) \why{left limit \eqref{eq:toolkit-left}}\\
&=F(\theta) \why{\(F\) is continuous}\\
\Pr(R=\theta)&=F(\theta)-F(\theta^-)=0 \why{the jump at \(\theta\); \(F\) has no jump}
\end{align*}
So the probability that the fixed value \(\theta\) beats \(R\) is \(F(\theta)\), and a tie has probability zero.
This is why the paper assumes a continuous \(F\): acquisition-value ties then have probability zero
(Section~2.2 of the paper).

\subsection{The probability that a prepared challenger wins}

The paper writes the probability that a prepared class-\(j\) challenger wins against the incumbent distribution
\(F\) as (Section~2.4 of the paper)
\begin{equation}\label{eq:toolkit-win}
W_j(F)=\E_j\bigl[F(\theta)\,\1\{\theta\ge p\}\bigr].
\end{equation}
The formula combines the tools above. The highest admissible bid wins, and a bid equal to the reserve is
admissible (Section~2.2 of the paper).

So a prepared challenger with the value \(\theta\) wins in two cases. If
\(R<p\le\theta\), it is the only admissible bidder. If \(p\le R<\theta\), it outbids the incumbent. Together, it
wins exactly when \(\theta\ge p\) and \(R<\theta\). Chapter~\ref{ch:auction} reads the auction rule in detail.

\begin{derivation}[the win probability]
Given the value \(\theta=y\):
\begin{align*}
\Pr(\text{win}\mid\theta=y)&=\1\{y\ge p\}\Pr(R<y) \why{\(R\) is independent of \(\theta\)}\\
&=\1\{y\ge p\}\,F(y) \why{\(F\) is continuous}
\end{align*}
Average over the class distribution \(G_j\), as in the law of iterated expectations
(Section~\ref{sec:toolkit-expectation}). The result is \(W_j(F)=\E_j[F(\theta)\1\{\theta\ge p\}]\).
\end{derivation}

With the uniform incumbent on \([0,r]\) and the example's point masses, where \(p<\ell<r<h\):
\begin{align*}
W_H&=F(h)\,\1\{h\ge p\}=1 \why{\(h>r\), so \(F(h)=1\); and \(h>p\)}\\
W_L&=F(\ell)\,\1\{\ell\ge p\}=\frac{\ell}{r} \why{\(p<\ell<r\), so \(F(\ell)=\ell/r\)}
\end{align*}
Class \(H\) wins for sure: in the paper's words, class \(H\) outbids the incumbent, \(W_H(F)=1\). Class \(L\) wins
less often against a stronger incumbent, because \(\ell/r\) falls as \(r\) rises.

\subsection{An illustration with the minimum}

The minimum formula gives a short route to the target's proceeds. Under the class ranking every challenger value
meets the reserve, and the target receives \(\max\{p,\min(R,\theta)\}\) after the challenger prepares
(Section~3.1 of the paper). Write \(Z=\min(R,\theta)\ge0\), where \(R\) and \(\theta\) are independent given the
class.
\begin{align*}
\Pr(Z>y)&=[1-F(y)]\,[1-G_j(y)] \why{minimum formula \eqref{eq:toolkit-mintail}}\\
t_j=\E_j\bigl[\max\{p,Z\}\bigr]&=p+\int_p^\infty\Pr(Z>y)\,dy \why{floor formula \eqref{eq:toolkit-max}}\\
&=p+\int_p^\infty[1-F(y)]\,[1-G_j(y)]\,dy \why{substitute the first line}
\end{align*}
The integral adds up the slices above the reserve that both bidders pass, since the target collects the second
bid. The result agrees with the paper's formula for \(t_j\) in Appendix~A.1 under the class ranking. There, the
two terms \(t_0+pF(p)[1-G_j(p^-)]\) reduce to \(p[1-F(p)]+pF(p)=p\), because \(G_j(p^-)=0\).
Chapter~\ref{ch:auction} derives the formula from the auction rule.

\begin{watchout}
\begin{itemize}
\item The product formulas need independence. In the paper, the challenger's class and value, the incumbent's
value, the cost and the draws of the order slot are mutually independent (Section~2.2 of the paper).
\item The win probability \(F(\theta)\) needs a continuous \(F\). With an atom of \(F\) at \(\theta\), the tie
\(R=\theta\) has positive probability, and the auction would need a rule to break it.
\end{itemize}
\end{watchout}

\begin{checkyourself}\label{q:toolkit-maxmin}
Suppose, against the paper's assumption, that the incumbent's value has an atom of probability \(a>0\) at a value
\(v\). A challenger with the value \(\theta=v\) bids against it.
(a)~Find \(\Pr(R<\theta)\) and \(\Pr(R\le\theta)\) in terms of \(F(v)\) and \(a\).
(b)~Why does the win probability then differ from \(F(\theta)\)?
\end{checkyourself}

% =====================================================================
\section{First-order stochastic dominance}\label{sec:toolkit-fosd}

\begin{intuition}[higher at every level]
One distribution dominates another when it makes high values more likely at every level. In terms of
distribution functions, the dominant distribution has a smaller \(G\) at every \(y\): it is less likely to be at
most \(y\). Then every increasing function of the value has a weakly larger average. For the incumbent, a stronger
incumbent is less likely to have a value below any level, so a challenger at a given value beats it less often.
\end{intuition}

\subsection{The definition}

A distribution \(\tilde G\) dominates a distribution \(G\) in the first order when
\begin{equation}\label{eq:toolkit-fosd}
\tilde G(y)\le G(y)\qquad\text{for every }y.
\end{equation}
Subtract both sides from one: \(1-\tilde G(y)\ge1-G(y)\). At every level, the dominant distribution passes the
level with at least the same probability. The paper uses the definition in three places.
\begin{itemize}
\item The class ranking \(G_H\le G_L\): class \(H\) dominates class \(L\) (Section~2.2 of the paper).
\item A strengthening \(\tilde F\le F\) of the incumbent (Proposition~1).
\item A path \(\{F_r\}\): \(F_{r'}\le F_r\) whenever \(r'>r\), so a larger \(r\) means a stronger incumbent
(Section~2.2 of the paper).
\end{itemize}

\subsection{Dominance raises the average of every increasing function}

\begin{keyidea}
Let \(X\ge0\), let \(\tilde G\le G\), and let \(\varphi\) be differentiable with \(\varphi'\ge0\), with finite expectations
under both laws. Then \(\E_{\tilde G}[\varphi(X)]\ge\E_G[\varphi(X)]\).
\end{keyidea}

This is the direction that the paper uses. The proof has two steps. The first step extends the tail formula to a
function \(\varphi\).

\begin{derivation}[the tail formula for an increasing function]
\begin{align}
\varphi(X)&=\varphi(0)+\int_0^X\varphi'(y)\,dy \why{fundamental theorem of calculus}\notag\\
&=\varphi(0)+\int_0^\infty \varphi'(y)\,\1\{y<X\}\,dy \why{the indicator cuts the integral at \(X\)}\notag\\
\E[\varphi(X)]&=\varphi(0)+\int_0^\infty \varphi'(y)\,\E\bigl[\1\{y<X\}\bigr]\,dy \why{linearity; swap, integrand \(\ge0\)}\notag\\
&=\varphi(0)+\int_0^\infty \varphi'(y)\,[1-G(y)]\,dy \whyinline{expectation of an indicator}\label{eq:toolkit-tailu}
\end{align}
\end{derivation}

The integral adds up, level by level, the gain \(\varphi'(y)\) from one more slice of value, weighted by the probability
that the value reaches the slice. With \(\varphi(x)=x\), the gain is one at every level, and \eqref{eq:toolkit-tailu}
is the tail formula~\eqref{eq:toolkit-tail}.

\begin{derivation}[the comparison]
\begin{align*}
&\E_{\tilde G}[\varphi(X)]-\E_G[\varphi(X)]\\
&\qquad=\int_0^\infty \varphi'(y)[1-\tilde G(y)]\,dy-\int_0^\infty \varphi'(y)[1-G(y)]\,dy \why{\eqref{eq:toolkit-tailu}; \(\varphi(0)\) cancels}\\
&\qquad=\int_0^\infty \varphi'(y)\,[G(y)-\tilde G(y)]\,dy \why{combine the integrals}\\
&\qquad\ge0 \why{\(\varphi'\ge0\) and \(G-\tilde G\ge0\)}
\end{align*}
\end{derivation}

In words: at each level, the dominant law adds the extra probability \(G(y)-\tilde G(y)\) of a pass, and each pass
is worth \(\varphi'(y)\ge0\). The finite expectations are no problem in the paper. The incumbent's values lie in
\([0,\bar r]\), and the class distributions have finite means (Section~2.2 of the paper).

The other direction holds too, if step functions count as increasing functions. The step function
\(\1\{x>y\}\) is nondecreasing in \(x\), and its expectation is \(1-G(y)\). If every nondecreasing function has a
weakly larger expectation under \(\tilde G\), then \(1-\tilde G(y)\ge1-G(y)\), which is \(\tilde G(y)\le G(y)\).

\subsection{The class difference \(D\)}

The paper writes \(D(y)=G_L(y)-G_H(y)\) (Section~2.2 of the paper). Rewrite it with tails:
\[
D(y)=[1-G_H(y)]-[1-G_L(y)].
\]
So \(D(y)\) is the extra probability that a class-\(H\) value passes the level \(y\). Under the class ranking,
\(D(y)\ge0\) at every \(y\). Apply the tail formula to each class and subtract:
\begin{align*}
\E_H[\theta]-\E_L[\theta]
&=\int_0^\infty[1-G_H(y)]\,dy-\int_0^\infty[1-G_L(y)]\,dy \why{tail formula \eqref{eq:toolkit-tail}, twice}\\
&=\int_0^\infty D(y)\,dy \why{combine the integrals}
\end{align*}
The mean gap between the classes is the sum of the slices at which class \(H\) passes more often.

For the example's point masses, with \(\ell<h\),
\[
D(y)=\1\{y\ge\ell\}-\1\{y\ge h\}=\1\{\ell\le y<h\}.
\]
Check the three regions. Below \(\ell\), both indicators are zero. From \(\ell\) up to \(h\), only the first is
one. From \(h\) on, both are one, and they cancel. The paper uses this form (Appendix~A.1). The mean gap is
\(\int_0^\infty\1\{\ell\le y<h\}\,dy=h-\ell\), as it should be.

\subsection{Uniform incumbents form a path}

Let \(r'>r>0\). The uniform law on \([0,r']\) dominates the uniform law on \([0,r]\). Compare
\(F_{r'}\) and \(F_r\) from \eqref{eq:toolkit-uniform} in four regions.
\begin{itemize}
\item \(y<0\): both are zero.
\item \(0\le y\le r\): \(F_{r'}(y)=y/r'\le y/r=F_r(y)\), because \(y\ge0\) and \(1/r'<1/r\). The inequality is
strict for \(y>0\).
\item \(r<y\le r'\): \(F_{r'}(y)=y/r'\le1=F_r(y)\), because \(y\le r'\).
\item \(y>r'\): both are one.
\end{itemize}
So \(F_{r'}\le F_r\) everywhere. In the paper's words, an increase in \(r\) strengthens the uniform incumbent, and
the strengths form a path (Section~4.2 of the paper).

\subsection{Why a stronger incumbent has a smaller \(F\)}

Read \(F(y)=\Pr(R\le y)\) from the challenger's side. By Section~\ref{sec:toolkit-maxmin}, a challenger with the
value \(\theta\) beats the incumbent with the probability \(F(\theta)\). A stronger incumbent is less likely to have
a value at most \(y\), at every level \(y\). So \(F\) is smaller at every level, and a challenger at any fixed value
wins less often. The paper states this in words: a stronger incumbent never raises the probability that a
prepared challenger wins (Section~3.1 of the paper).

Two distributions need not be ranked. As an illustration of our own, compare a point mass at \(a>0\) with the
uniform law on \([0,2a]\). Below
\(a\), the point mass has \(G=0\), and the uniform law has \(y/(2a)>0\) for \(y>0\). From \(a\) up to \(2a\), the
point mass has \(G=1\), and the uniform law has \(y/(2a)<1\). The two functions cross, so neither law dominates the
other.

\begin{watchout}
\begin{itemize}
\item Keep the direction straight. Dominance means a smaller distribution function. A strengthening of the
incumbent lowers \(F\).
\item Dominance is a comparison at every level. Two laws that cross are not ranked, and Proposition~1 does not
apply to them.
\item A strengthening can leave \(F\) unchanged on some levels, as the uniform laws do above \(r'\). So its effects
are weak in general. Proposition~1 states the integral conditions that make them strict.
\end{itemize}
\end{watchout}

\begin{checkyourself}\label{q:toolkit-fosd}
Let \(\tilde F\le F\) be two continuous incumbent distributions, and fix a challenger value \(\theta\ge p\).
(a)~Show that the challenger's expected gross profit \(\int_p^\theta F(y)\,dy\) of
Section~\ref{sec:toolkit-tail} is weakly smaller under \(\tilde F\).
(b)~Say in words why, with the slices picture.
\end{checkyourself}

% =====================================================================
\section{Games, equilibrium and deviations}\label{sec:toolkit-games}

\begin{intuition}[no one gains alone]
An equilibrium is a list of plans, one for each player, in which no player can gain by a change of its own plan
alone. Each player takes the plans of the others as given. Beliefs come from the plans by Bayes' rule. A mixed plan
is a lottery over actions, and a player uses a lottery only when the actions in it pay the same.
\end{intuition}

\subsection{Players, strategies and best responses}

\begin{itemize}
\item A \emph{player} is a decision maker. In the paper the players are the informed investor, the challenger
and the bidders in the auction. The market makers set the price equal to the expected value of a share given
the order, which the paper calls competitive prices.
\item A \emph{strategy} is a complete plan: an action for each situation that the player can face. The
investor's strategy gives, for each class \(j\), a distribution \(\sigma_j\) over the orders \(\{-1,0,1\}\). The
challenger's strategy says, for each price and each cost, whether it prepares (Section~2.4 of the paper).
\item A \emph{best response} is a strategy that gives the player the highest expected payoff, given the
strategies of the others.
\item A \emph{Nash equilibrium} is a list of strategies, one for each player, in which each strategy is a best
response to the others.
\end{itemize}
The paper's equilibrium (Definition~A.1) also fixes the beliefs and the prices, in six conditions.
Chapter~\ref{ch:model} reads them one by one.

\subsection{Mixed strategies and the indifference rule}

A mixed strategy puts the probability \(\sigma(a)\) on each action \(a\). Hold the strategies of the others fixed,
and write \(U(a)\) for the expected payoff of the action \(a\). By linearity, the mix pays the average
\(\sum_a\sigma(a)U(a)\).

\begin{keyidea}
A mix is a best response exactly when every action that it uses, with \(\sigma(a)>0\), has the highest payoff. In
particular, a player that mixes between two actions is indifferent between them.
\end{keyidea}

\begin{derivation}[the indifference rule]
Write \(U^*=\max_aU(a)\) for the highest payoff.
\begin{align*}
\sum_a\sigma(a)U(a)
&=\sum_a\sigma(a)U^*-\sum_a\sigma(a)\,[U^*-U(a)] \why{add and subtract \(\sum_a\sigma(a)U^*\)}\\
&=U^*-\sum_a\sigma(a)\,[U^*-U(a)] \why{the weights add up to one}
\end{align*}
Each term \(\sigma(a)[U^*-U(a)]\) is nonnegative. So the mix pays at most \(U^*\). It pays exactly \(U^*\) when
every term is zero, that is, when each action with \(\sigma(a)>0\) has \(U(a)=U^*\).
\end{derivation}

The paper's condition (E5) is this rule for the investor: \(\sigma_j(q)>0\) only if the order \(q\) maximizes
\(U_j(q)\) (Definition~A.1).

\subsection{Weakly dominant strategies}

A strategy is weakly dominant when it pays at least as much as any other strategy of the player. This must hold
against every possible choice of the others. The word ``weakly'' allows ties: against some choices of the others, another
strategy can pay the same. A weakly dominant strategy is a best response whatever the others do, so the player
needs no belief about them. In the paper's second-price auction, a bid equal to one's value is weakly dominant for
a bidder whose value meets the reserve (Section~2.2 of the paper). Chapter~\ref{ch:auction} proves it.

\subsection{Beliefs from strategies}

In an equilibrium each player knows the strategies of the others, but not their private information. So an
observer reads an action through the strategy. Suppose that an informed player sends a signal \(s\) with the
probability \(\sigma_H\) in class \(H\) and \(\sigma_L\) in class \(L\). Bayes' rule~\eqref{eq:toolkit-bayes} with
\(f_j(s)=\sigma_j\) gives the observer's posterior after \(s\):
\[
\mu(s)=\frac{\pi\sigma_H}{\pi\sigma_H+(1-\pi)\sigma_L}.
\]
Two cases show the range.
\begin{itemize}
\item If \(\sigma_H=\sigma_L>0\), then \(\mu(s)=\pi\). A plan that treats both classes alike reveals nothing.
\item If \(\sigma_H=1\) and \(\sigma_L=0\), then \(\mu(s)=1\). A plan that separates the classes reveals the
class.
\end{itemize}
The belief depends on the strategy that the observer expects. In the paper, a liquidity trader can place every
order, so the order that market makers see is a noisy signal. The posterior after a purchase or a sale therefore
stays inside the bounds of Section~\ref{sec:toolkit-odds}. The posterior after no order stays inside them when
the probability \(\beta\) of a liquidity purchase is at most \(1/3\). Chapter~\ref{ch:beliefs} shows how.

\subsection{Unilateral deviations}

To test an equilibrium, change the strategy of one player and hold everything else fixed. ``Everything else''
includes the strategies of the others, and also the prices and the entry probabilities that come from them. The
reason: the others do not see the change. They see only the action, and they read it through the strategy of the
candidate equilibrium. In the paper, the investor's deviation payoff is \(U_j(q)\) at the candidate's price
\(P(q)\) and entry \(e(q)\) (Section~2.4 of the paper; Appendix~A.2).

\begin{analogy}[a posted price list]
A shopper changes what goes into the basket, but the price list of the shop stays the same. The shop does not
know that this shopper changed plans. Where the analogy breaks: in the model the price comes after the order. It
comes from a rule, the candidate's price function. A deviation cannot change the rule, because market makers cannot
tell an order after a deviation from an order that they expect.
\end{analogy}

\subsection{Several equilibria and tie rules}

More than one list of strategies can pass the test. This can happen when a player's payoff depends on how the
others read its action, and the way they read it depends on the strategy. In the paper, a quiet market and an informative
market can both be equilibria in the same economy (Chapter~\ref{ch:coexistence}).

At an indifference, the definition must say what the indifferent player does. The paper's default tie rules are
these: an indifferent investor may mix over its best orders, and an indifferent high-cost challenger prepares. The
alternatives are that an indifferent investor abstains, and that an indifferent high-cost challenger may prepare
with any probability (Section~2.4 of the paper). Two rules for each player give four tie rules in total. Each
result states its tie rule. Proposition~2 holds under all four, and Propositions~3 and~4 use the default tie
rules.

\subsection{Comparisons across economies}

A comparison of two economies, such as a weak and a strong incumbent, solves the equilibrium again in each economy
(Section~2.4 of the paper). Prices, beliefs and entry can all differ between the two economies. A deviation is a
different exercise: it holds prices and entry fixed inside one economy.

\begin{watchout}
\begin{itemize}
\item Keep a deviation and a comparison apart. A deviation faces fixed prices and entry. A comparison across
incumbent strengths solves the equilibrium again in each economy.
\item The indifference rule applies to the actions that a mix uses. An action outside the mix may pay less.
\item State the tie rule with each result. A conclusion can depend on what an indifferent player does.
\end{itemize}
\end{watchout}

\begin{checkyourself}\label{q:toolkit-games}
In a candidate equilibrium, the investor in class \(H\) mixes between a purchase (\(q=1\)) and no order
(\(q=0\)).
(a)~What must hold for \(U_H(1)\), \(U_H(0)\) and \(U_H(-1)\)?
(b)~Suppose instead that \(U_H(1)>U_H(0)\). Why does the candidate fail?
(c)~When the investor tests a deviation to a pure purchase, which price and which entry probability does it face?
\end{checkyourself}

% =====================================================================
\section{Monotone functions, thresholds and strict inequalities}\label{sec:toolkit-monotone}

\begin{intuition}[one threshold, and room to spare]
Many arguments in the paper compare one quantity with a level. When the quantity moves in one direction only, it
crosses the level at most once, and one threshold splits the cases. A strict inequality has room to spare, so a
small change of the parameters cannot break it.
\end{intuition}

\subsection{Monotone functions}

A function \(\varphi\) is nondecreasing when \(x<x'\) implies \(\varphi(x)\le\varphi(x')\). It is strictly
increasing when \(x<x'\) implies \(\varphi(x)<\varphi(x')\). Nonincreasing and strictly decreasing functions are
the mirror cases. This chapter met three examples:
\begin{itemize}
\item the map \(o\mapsto o/(1+o)\) from odds to a probability is strictly increasing
(Section~\ref{sec:toolkit-odds});
\item every distribution function is nondecreasing (Section~\ref{sec:toolkit-cdf});
\item \(B(\mu)=g_L+\mu(g_H-g_L)\) is nondecreasing when \(g_H\ge g_L\) (Section~\ref{sec:toolkit-expectation}).
\end{itemize}

A monotone function turns a bound on its input into a bound on its output. The paper's exclusion bound is an
instance (Section~3.2 of the paper; Corollary~A.1). Suppose that \(g_H\ge g_L\), and that every belief \(\mu\) that
the challenger can hold satisfies \(\mu\le M(\pi)\). In the paper this holds after every order and at every price
when \(\beta\le1/3\), and Corollary~A.1 states this domain (Chapter~\ref{ch:beliefs}). If \(B(M(\pi))<c_H\), then at
every such belief
\begin{align*}
B(\mu)&\le B(M(\pi)) \why{\(B\) is nondecreasing and \(\mu\le M(\pi)\)}\\
&<c_H \why{assumption}
\end{align*}
So gross profit falls short of the high cost at every belief, and no price recruits a high-cost challenger.

\subsection{The positive part and the indicator}

Recall the positive part of a number \(x\), \((x)_+=\max\{x,0\}\) (Section~\ref{sec:toolkit-tail}). Its rules:
\begin{itemize}
\item \((x)_+\ge0\) and \((x)_+\ge x\);
\item \((x)_+=x\,\1\{x>0\}\);
\item \((x)_+\) is nondecreasing in \(x\);
\item \(\max\{a,b\}=b+(a-b)_+\) and \(\min\{a,b\}=a-(a-b)_+\).
\end{itemize}
Check the last rule in two cases. If \(a\ge b\), then \(\max\{a,b\}=a=b+(a-b)\) and \(\min\{a,b\}=b=a-(a-b)\). If
\(a<b\), then \((a-b)_+=0\), so the right sides are \(b\) and \(a\), the maximum and the minimum. In the paper, a
prepared challenger keeps \((\theta-\max\{p,R\})_+\) (Section~3.1 of the paper).

The indicator \(\1\{A\}\) equals one when \(A\) holds and zero otherwise. Its rules:
\begin{itemize}
\item \(\1\{A\}\,\1\{B\}=\1\{A\text{ and }B\}\), because the product is one only when both factors are one;
\item \(\1\{\text{not }A\}=1-\1\{A\}\);
\item \(\E[\1\{A\}]=\Pr(A)\), by \eqref{eq:toolkit-indicator}.
\end{itemize}
These rules give the point-mass difference of Section~\ref{sec:toolkit-fosd} in one more way, for \(\ell<h\):
\begin{align*}
\1\{\ell\le y<h\}
&=\1\{y\ge\ell\}\,\bigl(1-\1\{y\ge h\}\bigr) \why{``\(y\ge\ell\) and not \(y\ge h\)''}\\
&=\1\{y\ge\ell\}-\1\{y\ge\ell\}\,\1\{y\ge h\} \why{expand}\\
&=\1\{y\ge\ell\}-\1\{y\ge h\} \why{\(y\ge h\) implies \(y\ge\ell\)}
\end{align*}

\subsection{A straight line crosses a level once}

Take \(B(\mu)=g_L+\mu(g_H-g_L)\) from \eqref{eq:toolkit-B}, and suppose that \(g_H>g_L\).

\begin{derivation}[the preparation threshold]
\begin{align}
B(\mu)\ge c_H
&\iff g_L+\mu(g_H-g_L)\ge c_H \why{formula \eqref{eq:toolkit-B}}\notag\\
&\iff \mu(g_H-g_L)\ge c_H-g_L \why{subtract \(g_L\) from both sides}\notag\\
&\iff \mu\ge\frac{c_H-g_L}{g_H-g_L}=\tau \whyinline{divide by \(g_H-g_L>0\); the direction stays}\label{eq:toolkit-tau}
\end{align}
\end{derivation}

This is the paper's preparation threshold \(\tau\) (Section~3.2 of the paper): when \(g_H>g_L\), \(B(\mu)\ge c_H\)
if and only if \(\mu\ge\tau\). The negation of both sides gives \(B(\mu)<c_H\) if and only if \(\mu<\tau\). So the
exclusion condition \(B(M(\pi))<c_H\) says that the top posterior falls short of the threshold, \(M(\pi)<\tau\).

The sign of the slope matters.
\begin{itemize}
\item If \(g_H<g_L\), the last step divides by a negative number, and the inequality turns around:
\(B(\mu)\ge c_H\) if and only if \(\mu\le\tau\). This case does not arise under the class ranking, which gives
\(g_H\ge g_L\) (Appendix~A.1).
\item If \(g_H=g_L\), then \(B(\mu)=g_L\) at every belief. Either every belief passes, when \(g_L\ge c_H\), or none
passes, when \(g_L<c_H\). The threshold is not defined, and news about the class cannot change the challenger's
choice.
\end{itemize}
Check yourself~\ref{q:toolkit-monotone} asks where \(\tau\) lies relative to zero and one.

\subsection{Weak and strict inequalities}

The paper tracks the difference between \(\le\) and \(<\) with care, for three reasons.
\begin{itemize}
\item \emph{Ties.} Condition (E3) fixes the high-cost challenger's choice when \(B>c_H\) or \(B<c_H\). At
\(B=c_H\), a tie rule decides. Under the default rule, the indifferent challenger prepares, so it prepares exactly
when \(\mu\ge\tau\), a weak inequality (Lemma~A.2(d)).
\item \emph{Weak conclusions.} Proposition~1 says that a strengthening weakly lowers \(g_j\) and weakly raises the
spread. A strict change needs the integral conditions that the proposition states.
\item \emph{Strict hypotheses.} The hypotheses (H1)--(H3) of Proposition~2 are strict inequalities, and every
comparison behind the result is strict. So the outcome is the same under all four tie rules (Section~4.1 of the
paper).
\end{itemize}

\subsection{The intermediate value theorem}

Let \(\varphi\) be continuous on an interval \([a,b]\): its graph has no jumps. If \(\varphi(a)<k<\varphi(b)\), or
\(\varphi(a)>k>\varphi(b)\), then \(\varphi(x)=k\) at some \(x\) between \(a\) and \(b\). A graph without jumps
cannot get from below the level to above it without a point on the level. If \(\varphi\) is also strictly
monotone, this point is unique, because a strictly monotone function never takes the same value twice.

As an illustration, let \(g_L<c_H<g_H\). The function \(B\) is continuous on \([0,1]\), with \(B(0)=g_L<c_H\) and
\(B(1)=g_H>c_H\). The theorem gives a belief in \((0,1)\) at which \(B=c_H\). Since \(B\) is strictly increasing,
there is only one such belief, and formula~\eqref{eq:toolkit-tau} names it: \(\tau\). The paper uses the theorem
in the same way for the probability with which the investor mixes (Appendix~A.4; Chapter~\ref{ch:coexistence}).

\subsection{An open set of parameters}

\begin{keyidea}
Let \(f\) and \(g\) be continuous functions of a list of parameters \(x\), and let \(f(x_0)>g(x_0)\). Then
\(f(x)>g(x)\) for every \(x\) near enough to \(x_0\).
\end{keyidea}

\begin{derivation}[strict inequalities survive small changes]
Write \(d=f(x_0)-g(x_0)>0\) for the gap. By continuity, every \(x\) near enough to \(x_0\) has
\(|f(x)-f(x_0)|<d/2\) and \(|g(x)-g(x_0)|<d/2\). For such \(x\):
\begin{align*}
f(x)&>f(x_0)-\tfrac d2 \why{\(f\) moves by less than \(d/2\)}\\
&=g(x_0)+\tfrac d2 \why{\(f(x_0)=g(x_0)+d\)}\\
&>g(x) \why{\(g\) moves by less than \(d/2\)}
\end{align*}
\end{derivation}

A finite list of strict inequalities also survives: take the smallest of the neighborhoods. So the set of
parameters at which the list holds is open: each point in it has a small neighborhood inside it. A weak
inequality has no room to spare. At an equality \(f(x_0)=g(x_0)\), a small change can push \(f\) below \(g\).

The paper uses this argument for its main result. Fix the example's classes and reserve, and let the incumbents
be uniform. Then the hypotheses (H1)--(H3) are strict inequalities between continuous functions of
\((r_0,r_1,r_2,c_H,k,\rho,\pi,\lambda,\beta)\). So they hold on an open set of these parameters that contains the
example (Section~4.2 of the paper; Appendix~A.3). In words, the result is not a
knife-edge case: small changes of the parameters keep it.

\begin{watchout}
\begin{itemize}
\item Division by \(g_H-g_L\) keeps the direction of an inequality only when \(g_H-g_L>0\).
\item ``Open'' says that small changes keep the result. It says nothing about the size of the set.
\item An open set needs both strictness and continuity. A weak inequality, or a function with a jump, can fail
after a small change.
\end{itemize}
\end{watchout}

\begin{checkyourself}\label{q:toolkit-monotone}
Let \(g_H>g_L\).
(a)~Describe the set of beliefs \(\mu\in[0,1]\) with \(B(\mu)\ge c_H\) in three cases: \(c_H\le g_L\);
\(g_L<c_H\le g_H\); and \(c_H>g_H\).
(b)~Suppose that every belief that the challenger can hold lies in \([m(\pi),M(\pi)]\). Which condition on \(\tau\)
rules out high-cost preparation at every such belief?
\end{checkyourself}

% =====================================================================
\begin{chaptersummary}
Bayes' rule updates the belief about the class, and in odds form it multiplies the prior odds by the likelihood
ratio. A bound \(1+u\) on every likelihood ratio keeps every posterior, and every average of posteriors, between
\(m(\pi)\) and \(M(\pi)\). A distribution function and its left limit give the probabilities of ``at most'' and
``below'', and the tail formula turns them into expectations, slice by slice. Every expectation at a belief is a
straight line in the belief when the signal is independent of the value given the class, so
\(B(\mu)=g_L+\mu(g_H-g_L)\). A stronger incumbent has a smaller distribution function at every level, so a
challenger at any fixed value wins less often. An equilibrium is a list of best responses with consistent
beliefs: a deviation holds prices and entry fixed, and a comparison of economies solves each economy again. A
straight line with a positive slope crosses a level once, at the threshold \(\tau\), and strict inequalities between
continuous functions survive small changes of the parameters.
\end{chaptersummary}

% =====================================================================
\answersection

\answer{\ref{q:toolkit-bayes}}{%
(a)~Put \(f_H(s)=f_L(s)=f\) into formula~\eqref{eq:toolkit-bayes}.
\begin{align*}
\mu(s)&=\frac{\pi f}{\pi f+(1-\pi)f} \why{substitute}\\
&=\frac{\pi f}{f} \why{\(\pi f+(1-\pi)f=f\)}\\
&=\pi \why{cancel \(f>0\)}
\end{align*}
A signal that is equally likely in both classes carries no news about the class.\par
(b)~Put \(f_L(s)=0\): \(\mu(s)=\pi f_H(s)/[\pi f_H(s)+0]=1\). Only class \(H\) can send this signal.\par
(c)~As \(\pi\) approaches one, the term \((1-\pi)f_L(s)\) in the denominator approaches zero, while \(\pi f_H(s)\)
approaches \(f_H(s)>0\). So the posterior approaches one. A strong prior stays strong after one signal of bounded
strength.}

\answer{\ref{q:toolkit-odds}}{%
(a)~Start with the top posterior.
\begin{align*}
M(\pi)>\pi
&\iff\frac{1+u}{\pi(1+u)+1-\pi}>1 \why{divide by \(\pi>0\)}\\
&\iff 1+u>\pi(1+u)+1-\pi \why{multiply by the positive denominator}\\
&\iff(1-\pi)(1+u)>1-\pi \why{subtract \(\pi(1+u)\) from both sides}\\
&\iff u>0 \why{divide by \(1-\pi>0\), subtract one}
\end{align*}
For the bottom posterior, \(m(\pi)<\pi\) is the same as \(1/[\pi+(1-\pi)(1+u)]<1\), after division by \(\pi>0\).
This is \(\pi+(1-\pi)(1+u)>1\), that is \((1-\pi)(1+u)>1-\pi\), that is \(u>0\).\par
(b)~Put \(\pi=1/2\).
\begin{align*}
M(1/2)&=\frac{\tfrac12(1+u)}{\tfrac12(1+u)+\tfrac12}=\frac{1+u}{2+u} \why{multiply above and below by \(2\)}\\
m(1/2)&=\frac{\tfrac12}{\tfrac12+\tfrac12(1+u)}=\frac{1}{2+u} \why{multiply above and below by \(2\)}
\end{align*}
The sum is \((1+u+1)/(2+u)=1\). This also follows from the mirror identity~\eqref{eq:toolkit-mirror} at
\(\pi=1/2\), where \(1-\pi=\pi\).}

\answer{\ref{q:toolkit-cdf}}{%
(a)~All the probability of class \(L\) sits at \(\ell=p\). No value lies strictly below \(p\), so
\(G_L(p^-)=\Pr(\theta<\ell)=0\). Every value is at most \(p\), so \(G_L(p)=\Pr(\theta\le\ell)=1\).\par
(b)~Yes. A bid equal to the reserve is admissible, and by \eqref{eq:toolkit-geq},
\(\Pr(\theta\ge p)=1-G_L(p^-)=1\). The number \(G_L(p)=1\) would wrongly suggest that every value fails, because it
also counts the values equal to \(p\).\par
(c)~The class ranking restricts the left limit, \(G_L(p^-)=0\), and the restriction holds here. The example itself
keeps a gap, \(p<\ell\) (Section~4.2 of the paper).}

\answer{\ref{q:toolkit-expectation}}{%
(a)~Substitute \(w_j=t_j-t_0\) and collect the terms.
\begin{align*}
w(\mu)&=(1-\mu)(t_L-t_0)+\mu(t_H-t_0) \why{definitions}\\
&=(t_L-t_0)-\mu(t_L-t_0)+\mu(t_H-t_0) \why{expand the first product}\\
&=t_L-t_0+\mu\,[(t_H-t_0)-(t_L-t_0)] \why{collect the terms in \(\mu\)}\\
&=t_L-t_0+\mu\,(t_H-t_L) \why{\(t_0\) cancels inside the bracket}\\
&=t_L-t_0+\DT\,\mu \why{\(\DT=t_H-t_L\)}
\end{align*}\par
(b)~At \(\mu=0\), \(w(0)=t_L-t_0=w_L\). At \(\mu=1\), \(w(1)=t_L-t_0+t_H-t_L=w_H\). The slope is the target-payoff
spread \(\DT\). The function is a straight line for the same reason as \(B(\mu)\): the belief only moves the
weights on the two classes.}

\answer{\ref{q:toolkit-tail}}{%
(a)~On \([p,\ell]\) every level lies below \(r\), so \(F(y)=y/r\) there.
\begin{align*}
g(\ell)=\int_p^\ell F(y)\,dy&=\int_p^\ell\frac yr\,dy \why{\(\ell<r\), so \(F(y)=y/r\) on \([p,\ell]\)}\\
&=\Bigl[\frac{y^2}{2r}\Bigr]_p^\ell \why{antiderivative of \(y/r\)}\\
&=\frac{\ell^2-p^2}{2r} \why{upper limit minus lower limit}
\end{align*}
This is the paper's \(g_L\) in the example (Section~4.2 of the paper).\par
(b)~The numerator \(\ell^2-p^2\) is positive and does not depend on \(r\), and the denominator \(2r\) rises with
\(r\). So \(g(\ell)\) falls as the incumbent strengthens, and it approaches zero when \(r\) grows without limit. In
slices: a stronger uniform incumbent lies below each level \(y\) less often, so the challenger keeps each slice less
often. As \(\ell\) approaches \(p\), the numerator approaches zero: a value just above the reserve leaves almost no
slices to keep. The formula needs \(r>\ell\). For \(r\le\ell\), the integrand is one on the part of \([p,\ell]\)
above \(r\), and the integral changes form.}

\answer{\ref{q:toolkit-maxmin}}{%
(a)~The atom makes \(F\) jump by \(a\) at \(v\), so \(F(v^-)=F(v)-a\). Then
\(\Pr(R<\theta)=F(v^-)=F(v)-a\), by \eqref{eq:toolkit-left}, and \(\Pr(R\le\theta)=F(v)\). The two differ by
\(\Pr(R=v)=a>0\).\par
(b)~With the probability \(a\), the two values are equal, and the two bids tie. The winner then depends on a rule
that breaks ties. So the win probability can lie anywhere between \(F(v)-a\) and \(F(v)\), and the formula
\(F(\theta)\) does not follow. A continuous \(F\) removes the problem: ties have probability zero.}

\answer{\ref{q:toolkit-fosd}}{%
(a)~Write the difference as one integral.
\begin{align*}
\int_p^\theta F(y)\,dy-\int_p^\theta\tilde F(y)\,dy
&=\int_p^\theta[F(y)-\tilde F(y)]\,dy \why{combine the integrals}\\
&\ge0 \why{\(F(y)-\tilde F(y)\ge0\) at every \(y\), and \(\theta\ge p\)}
\end{align*}
So the expected gross profit at the value \(\theta\) is weakly smaller under \(\tilde F\).\par
(b)~The challenger keeps the slice of value at the level \(y\) when the incumbent's value lies below \(y\). A
stronger incumbent lies below each level with a weakly smaller probability, so the challenger keeps each slice
weakly less often. The fall is strict when \(\tilde F<F\) on a set of levels between \(p\) and \(\theta\) of
positive length. Proposition~1(a) states the same comparison for the class averages \(g_j\).}

\answer{\ref{q:toolkit-games}}{%
(a)~By the indifference rule, the two orders in the mix pay the same, and no order pays more:
\(U_H(1)=U_H(0)\ge U_H(-1)\).\par
(b)~If \(U_H(1)>U_H(0)\), the mix pays \(\sigma_H(1)U_H(1)+\sigma_H(0)U_H(0)\), which is below \(U_H(1)\) because
\(\sigma_H(0)>0\). A pure purchase pays more, so the mix is not a best response. Condition (E5) fails, and the
candidate is not an equilibrium.\par
(c)~It faces the candidate's price after a purchase, \(P(1)\), and the candidate's entry after a purchase,
\(e(1)\). Market makers and the challenger do not see the change of plan, so they keep the price and the entry
rule of the candidate.}

\answer{\ref{q:toolkit-monotone}}{%
(a)~By \eqref{eq:toolkit-tau}, the set is \(\{\mu\in[0,1]:\mu\ge\tau\}\), with \(\tau=(c_H-g_L)/(g_H-g_L)\). The
denominator is positive, so the sign of \(\tau\) is the sign of \(c_H-g_L\), and \(\tau\le1\) exactly when
\(c_H-g_L\le g_H-g_L\), that is \(c_H\le g_H\).
\begin{itemize}
\item \(c_H\le g_L\): then \(\tau\le0\), and every belief in \([0,1]\) passes. Even a challenger sure of class
\(L\) covers the high cost.
\item \(g_L<c_H\le g_H\): then \(0<\tau\le1\), and the beliefs that pass form the interval \([\tau,1]\).
\item \(c_H>g_H\): then \(\tau>1\), and no belief in \([0,1]\) passes. Even a challenger sure of class \(H\) does
not cover the high cost.
\end{itemize}\par
(b)~High-cost preparation needs a belief \(\mu\ge\tau\). The largest belief that the challenger can hold is
\(M(\pi)\). So the condition \(M(\pi)<\tau\) rules out high-cost preparation at every belief. By
Section~\ref{sec:toolkit-monotone}, this is the same as \(B(M(\pi))<c_H\), the paper's exclusion condition.}
